\section{後期のテスト}
\label{lesson30}

\subsection{授業内容}
\subsubsection{科目の中でのこのコマの位置づけ}
後期で学んできたことについて，基本的な知識，用語の理解，心理統計の基本的な考え方の筋道などを等テストをする。
\subsection{授業情報}
\paragraph{コマの展開方法}
マークシート式の試験を行う。問題数は60-70問程度である。
\paragraph{諸注意}
\begin{itemize}
    \item 授業を行なっている教室で試験を実施する。着席場所は指定しないが，机の両端に座り，他の受験生と隣接しないようにする。
    \item 試験に際して資料等の持ち込みは認めない。
    \item 携帯電話，タブレット，PCなど通信機能を有する機器の使用も認めない。
    \item (関数）電卓がなくても計算ができるよう工夫された出題をしているため，電卓などの使用も認めない。
    \item 受験者数を確定するため，当日は出席管理システムへの登録を行なってもらう。
    \item 不正行為を見つけた場合は即時退出を求め，単位を与えない。
    \item 試験会場に20分以上遅刻したものは受験を認めない。
    \item 試験開始から30分以上経過すると，希望者には退出を許可する。
    \item 退出に当たっては，質問用紙と解答用紙の両方を提出すること。
    \item マークシートへの回答であるため，鉛筆の利用が望ましい。
    \item マークを全て塗りつぶしていないなど，読み取り上のエラーには対応しない。
    \item エクストラクレジットについての説明は別途行う。
\end{itemize}
